\documentclass[11pt, a4paper]{article}

\usepackage[utf8]{inputenc}
\usepackage[T1]{fontenc}
\usepackage{lmodern}
\usepackage[margin=2.4cm]{geometry}
\usepackage{amsmath, amssymb, amsthm}
\usepackage{booktabs}
\usepackage{enumitem}
\usepackage{hyperref}
\usepackage{xcolor}
\usepackage{microtype}
\usepackage{array}
\usepackage{longtable}
\usepackage{colortbl}
\usepackage{tabularx}

\hypersetup{
    colorlinks=true,
    linkcolor=black,
    citecolor=blue!60!black,
    urlcolor=blue!60!black
}

\definecolor{rowgreen}{RGB}{232,245,233}
\definecolor{rowred}{RGB}{253,237,237}
\definecolor{rowamber}{RGB}{255,248,225}

\newcommand{\M}{\mathbf{M}}
\newcommand{\R}{\mathbb{R}}

\newcommand{\rwin}{\cellcolor{rowgreen}}
\newcommand{\rloss}{\cellcolor{rowred}}
\newcommand{\rmixed}{\cellcolor{rowamber}}

\title{%
    \textbf{A Theory of Learning-Experiments:}\\[4pt]
    \Large What Bears Fruit, What Doesn't, and Why\\[2pt]
    \large Patterns from a Year of Orbital Response Network Research%
}
\author{%
    Sirsh Amarteifio\\
    {\small Percolation Labs}\\
    {\small \url{https://github.com/Percolation-Labs/orn}}%
}
\date{April 2026}

\begin{document}
\maketitle

\begin{abstract}
After a year of designing, training, and ablating architectures in the Orbital Response Network (ORN) programme, a pattern has emerged across roughly forty distinct experiments. Some recipes ,  shared coupling geometry, wide shared FFN, permutation-invariant structural probes, depth-distributed gates ,  reliably produce measurable gains. Others ,  polynomial operators of the coupling matrix, iterative sequential field recurrences, strict-orbit gauge constraints, 1D position encodings on 2D data ,  fail with equal reliability, often in ways that feel inevitable in hindsight. This paper does not report a new architecture. It steps back to ask: \emph{what distinguishes a promising experiment from an unpromising one before we run it?}

We cluster our experimental history into six motifs and propose a rulebook. The central empirical claim is that the productive experiments share a common structural property: they exploit an \textbf{invariant} (symmetry, shared substructure, order-independence) that the dominant training signal ,  next-token prediction ,  cannot itself see. The unproductive experiments share a dual property: they try to replace, at zero parameter cost, a capability that depth and parallelism make nearly free on the GPU. We then map these findings to two biological frameworks ,  hemispheric lateralisation (McGilchrist) and complementary learning systems (McClelland, O'Reilly, Kumaran) ,  and argue that both converge on the same selection principle. The paper closes with a short playbook for triaging future experiments against a checklist rather than vibes.
\end{abstract}

\tableofcontents

% ======================================================================
\section{Why a Theory of Experiments}

The ORN programme began with a single hypothesis: that per-layer query--key attention is largely redundant and can be memoised into a single shared metric $\M$. Validating that hypothesis required running a family of subsidiary experiments ,  variants of sharing, variants of coupling, variants of inference-time adaptation, variants of cross-modal transfer, variants of gauge-fixed operator surgery. Each of these had to be chosen from a much larger space of things we \emph{could} have tried. Some choices were right; many were wrong.

What changed over the year is not our theoretical taste but our \emph{priors}. We now know, before running an experiment, whether it is likely to produce a signal. The point of this paper is to make those priors explicit. Every architecture researcher accumulates them; few write them down; almost none reconcile them with biological precedent and hardware reality. Doing so matters because experiments are expensive. A single 3B-token run on an A100 is \$300 and a week. Cutting even a quarter of the unpromising candidates is worth a paper.

The method is empirical. We take the ORN programme's experimental log as a dataset, cluster the outcomes, and look for shared features. The result is not a proof ,  experimental programmes are too small for statistical claims ,  but a set of \emph{heuristics with a story}. Each heuristic has a mechanism-level justification and a biological or hardware correlate. Taken together they form a theory of what can be \emph{learned} in a gradient-based, GPU-parallel, next-token-predicting system.

% ======================================================================
\section{The Experimental Record}
\label{sec:record}

We list the principal experiments of the programme below, with outcome and one-line justification. Green rows are productive; red rows are negative results; amber rows are partial or conditional wins. This is not exhaustive ,  smaller sanity checks are omitted ,  but it covers every experiment that materially shaped the research direction.

\subsection{Table: What Was Tried}

\begin{small}
\begin{longtable}{p{4.2cm} p{1.5cm} p{8.6cm}}
\toprule
\textbf{Experiment} & \textbf{Outcome} & \textbf{Why (one-line)} \\
\midrule
\endhead
\multicolumn{3}{l}{\textit{Shared substructure}} \\
\midrule
\rwin SharedM (one $A,B$ for all layers) & works & 48$\times$ coupling compression, $<$1\% quality loss at matched params; the coupling \emph{rule} is invariant across layers. \\
\rwin Wide shared FFN (one $8\times d$ FFN for all layers) & works & One wide memory beats $L$ narrow copies at matched params; $-0.8\%$ val loss. \\
\rwin Asymmetric $\M = AB^\top$ vs symmetric $LL^\top$ & works & Language is directional; autoregressive tasks punish symmetric coupling. \\
\rloss ALBERT-style full sharing (Q, K, V, O, FFN) & fails & Shares too much, including the per-layer response; loses 2--4\% quality at matched params. \\
\rwin Frozen-M multimodal transfer (LoRA 0.83\%) & works & Coupling geometry generalises; audio adapts fast, image slower but real. \\
\midrule
\multicolumn{3}{l}{\textit{Depth, response, and corrections}} \\
\midrule
\rwin Perturbative context-gated corrections on residual stream & works & 0.449 nat ablation impact; gate anticorrelates with attention entropy ($r=-0.95$). \\
\rwin Per-layer response heads $W_V, W_O$ + per-layer FFN & works & Shared structure $+$ per-layer response is the right decomposition. \\
\rwin Test-time training (three-timescale) & works & Regime C (slow $\M$ drift, fast correction SGD): $-0.139$ nat at trained length. \\
\rmixed Deep R (second hemisphere, non-causal) & partial & Wins on relational tasks ($+94.8\%$ on dependency detection), loses on per-position tasks; useful but task-dependent. \\
\rloss Operator-space replacements for depth (MatExp, Spectral, Poly, S4, Deferred) & fail & All 5--9\% worse than 8-layer baseline; depth is labour that algebra cannot do in one shot. \\
\midrule
\multicolumn{3}{l}{\textit{Coupling as a functional class}} \\
\midrule
\rloss Polynomial of coupling $p(C) \cdot V$ for multi-hop & fails & Rank-compression of a normalised matrix; higher degrees vanish; gate stuck at 0.5. \\
\rloss Rank-1 collapse in spectral-normalised polynomial $\M$ & fails & Fast polynomial coefficients adapt to low-rank input, removing gradient pressure on $\M$. \\
\rloss Iterative field recurrence $G W G^\top$ per position & fails & $43\times$ slower; spectral version still $10\times$; no parallelism means GPUs punish it. \\
\rloss Abelian attention (sequential field cascade) & fails & Cascade length stuck at max; field never equilibrates; sequential semantics kills throughput. \\
\rwin Spectral rotation for coupling dynamics (eigendecomposition of $G$) & works & Structured aging; $|\lambda|$ range $[0, 0.51]$; phases input-dependent, eigenvalues not. \\
\midrule
\multicolumn{3}{l}{\textit{Gauge-LORN / structural steering}} \\
\midrule
\rloss Strict-orbit gauge (Path A K-sweep on $\M$) & fails & Cosines $\geq 0.9999$ across $K$; direction-uniform; nothing to latch onto. \\
\rloss FFN as gauge target & fails & High diagnostic z-score (4.41) but diffuse; no compact subspace. \\
\rloss Joint $\M$+FFN gauge & fails & Substrates do not synergise; $+2$ nats vs $\M$-only. \\
\rloss Response composition $R_l = W_v W_o$ as gauge target & fails & Highest z-score but collapsed to direction-uniform under learning. \\
\rwin Recoverable-orbit gauge (unit-norm projection on $\alpha$) v10 & works & 83\% category accuracy at $+0.06$ nat val; LDA basis in $\M$-space orthogonal to NTP gradient. \\
\rwin Permutation-invariant fingerprint (v17b) & works & 37 s train, genre accuracy 82\% on 7 classes; $S_n$-invariance is orthogonal to NTP's gradient. \\
\rwin Temporal-binding (source episodes, wiki vs OWT) & works & No category labels; co-occurrence alone suffices; outperforms label-supervised variants. \\
\midrule
\multicolumn{3}{l}{\textit{Hemispheric injection modes (R $\to$ L)}} \\
\midrule
\rwin Coupling bias injection (R biases L's attention logits) & works & Val 4.54, gate 86\%; structural template, not content. \\
\rwin Per-layer gated coupling (independent gate per deep layer) & works & Deep layers open gates to 0.999, shallow to 0.000; phase change at mid-depth. \\
\rloss Additive residual injection (R adds to L's stream) & fails & L closes gate to 0.16; content bias is ignorable. \\
\rmixed Cross-attention (L queries R) & partial & Gate 0.90 but val unchanged; L extracts what's easy, not what's structural. \\
\rmixed Value modulation (R gates L's values) & partial & Highest correlation with structure, weakest channel for prediction. \\
\midrule
\multicolumn{3}{l}{\textit{Multimodal / representation}} \\
\midrule
\rloss 1D RoPE on image patches & fails & R@1 at chance; 1D sequencing fights 2D spatial geometry. \\
\rloss CLS pooling with causal attention & fails & Position 0 self-attends only; blind to content; contrastive loss has no gradient. \\
\rwin Sub-linear multimodal rank growth & works & 3 modalities need rank 3 of $\M$, not 10; coupling geometry is partially shared. \\
\rwin Shared FFN owns 97\% of manifold, $\M$ owns 36\% & works & FFN is the world model, $\M$ is the compass; both shared components conspire. \\
\midrule
\multicolumn{3}{l}{\textit{Training recipe}} \\
\midrule
\rwin WSD (warmup-stable-decay) schedule & works & $-0.23$ nat in final 20\%; ARC-Easy jumps from 30.5\% to 46.2\% without weight change. \\
\rwin SDPA (fused attention) $+$ AMP $+$ grad checkpointing & works & Training throughput to 2.9K tok/s on A100; MFU 52\%. \\
\rloss Manual QKV attention with explicit masks & fails & Slow, leakage-prone on multimodal; SDPA fuses mask and softmax. \\
\bottomrule
\end{longtable}
\end{small}

% ======================================================================
\section{Clustering the Outcomes}
\label{sec:clusters}

Reading the table as a whole, six clusters emerge. Each cluster contains experiments of varying surface form that share a common mechanism.

\subsection{C1. Invariance beats redundancy}

SharedM, wide shared FFN, permutation-invariant fingerprints, temporal binding, frozen-M transfer ,  all win by identifying a quantity that is \emph{invariant} across the dimension the model is redundant over (layers, permutations, modalities, episodes). The standard transformer parameterises this invariant $L$ times; we parameterise it once and spend the saving on width or on a qualitatively different component.

The principle: \textbf{if a function is invariant under a transformation, express it as a single shared object; do not train $L$ copies and hope gradient descent ties them.}

This explains both sides of the ledger. ALBERT shares too aggressively ,  it assumes the \emph{response} is invariant, which is false, and pays for it. SharedM shares exactly the coupling rule, which is invariant. The distinction is empirical but principled: which component's Jacobian changes across layers, and which doesn't?

\subsection{C2. Orthogonality to the dominant gradient}

Every gauge-LORN probe that failed (v8, v9, Path A K-sweep, FFN target, response target, joint targets) had the same symptom: cluster cosines $\geq 0.9999$. The structural axis collapsed onto the direction the network already spent its capacity on ,  the next-token-prediction gradient. The probes that \emph{worked} (v10 unit-norm, v17b permutation-invariant, temporal binding) were those that found subspaces NTP's gradient cannot reach by construction.

The principle: \textbf{a training signal cannot discover structure that lies parallel to the signal it already optimises. To find structural axes, construct probes that live in the orthogonal complement of NTP's gradient flow.}

Unit-norm projection forces category encoding into direction, not magnitude; permutation averaging wipes out word order, leaving only the order-invariant residue (roughly 80\% of genre signal, empirically); temporal binding uses co-occurrence rather than category labels, which NTP cannot distinguish from local lexical context. All three exploit the same orthogonality trick.

\subsection{C3. Depth is labour, not ornament}

The operator-space cluster (MatExp, Spectral, Poly, S4, Deferred ORN, iterative field recurrence, abelian attention) tried to replace $L$ layers with one algebraic operation computed in closed form. Every variant underperformed. The diagnosis is not subtle: shallow layers are 96\% linear and accumulate; deep layers are 50\% nonlinear and reorganise. A single-pass algebra cannot provide both the accumulation and the re-coupling.

The principle: \textbf{depth is where representations are progressively enriched by the previous layer's nonlinear output. A closed-form operator-space substitute has no mechanism for re-coupling after the first nonlinearity.}

Crucially, this does not mean ``depth is mandatory for all problems.'' It means depth provides \emph{multi-level indirection}: couple, transform, recouple, transform. The experiment that replaces depth with algebra must also provide recoupling, which requires alternating linear and nonlinear operations ,  which is depth.

\subsection{C4. GPU-parallel beats sequential-algebraic}

The iterative recurrences (abelian attention, $G W G^\top$ per position) fail not only for representational reasons but because they saturate the accelerator's serial dependency chain. A $43\times$ slowdown means a $43\times$ ratio of unproductive floating-point operations per useful one. Even when the mathematics is correct, an architecture that cannot parallelise along sequence or head dimensions is not a viable candidate ,  it runs on the wrong hardware.

The principle: \textbf{any architecture whose per-step compute cannot be decomposed into a small number of large matmuls on the GPU is already dead, regardless of its mathematical merits.} Novel operations must either (a) reduce to GEMMs, (b) run at small enough scale to be dominated by overhead elsewhere, or (c) accept a constant-factor throughput hit that can be amortised against a representational win.

This is the PyTorch--GPU filter. It has eliminated an enormous fraction of our candidate designs before they reached evaluation, because what the GPU wants ,  wide, dense, batched, fused matmuls ,  is a strong constraint on what architectures are worth proposing.

\subsection{C5. Emergence beats prescription}

The per-layer gates in two-hemisphere coupling, the phase change at layer $\sim L/2$, the 0.85 saturation of the cooperative gate, the spectral collapse of the coupling operator to rank $\sim 70$, the ORN's 60$\times$ syntactic state separation ,  none of these were designed. Each emerged when we \emph{built a surface} that the optimiser could populate and let training do the work.

The dual failures are instructive: the polynomial-degree activation had its coefficients $c_2, c_3 \to 0$ ,  the optimiser rejected higher-degree structure when the task did not require it. The rank-1 collapse of spectral-normalised $\M$ happened because the response pathway could compensate, removing gradient pressure on $\M$.

The principle: \textbf{design the \emph{affordance}, not the \emph{behaviour}. Build a parameterisation that permits the target phenomenon, choose a loss and data regime that rewards it, and observe whether optimisation discovers it. Prescribing the behaviour directly (hand-chosen polynomial coefficients, fixed spectral shapes, designed operator templates) removes the optimiser's degrees of freedom and almost always fails.}

\subsection{C6. Diagnose before scaling}

The V3 multimodal run flew blind without modality-specific validation for ten days and cost \$30--50 plus four days of wall time on an experiment that was already broken (1D RoPE on patches, CLS+causal interaction). In every case where we caught a failure early, the reason was a cheap local probe: text generation on small checkpoints, per-modality loss curves, small-scale recovery tests.

The principle: \textbf{every experiment must ship with a cheap diagnostic that can detect failure within one epoch. Running a multi-day experiment without a same-day canary is not a parameter-efficient use of a bank account.}

The corollary: many failed experiments would have been killed in an hour if we had looked at the first-step gradients, the initial coupling pattern, the per-layer gate values, the first hundred generated tokens. Scaling without diagnostics is more expensive than it looks.

% ======================================================================
\section{A Rulebook for Triaging Experiments}
\label{sec:rulebook}

The clusters of Section~\ref{sec:clusters} turn into a checklist. Before committing compute to an architectural idea, we now ask:

\begin{enumerate}[label=\textbf{R\arabic*.}, leftmargin=*]
\item \textbf{What is the invariant?} Is the proposed component removing a redundancy (shared structure across layers, modalities, episodes)? If yes, proceed. If the proposal instead adds per-layer variation without a story, revisit.

\item \textbf{Where does the probe sit relative to the NTP gradient?} If the diagnostic or structural target lives in a subspace NTP can already optimise into directly, the result will collapse to cosine 1.0. Redesign the probe to be orthogonal by construction (permutation invariance, temporal binding, magnitude-normalisation, out-of-vocabulary supervision).

\item \textbf{Does the design respect depth as labour?} If it proposes to replace $L$ layers with a one-shot algebraic operation, reject unless it reinstates recoupling. Operator-space is not a substitute for alternating linear--nonlinear stacking.

\item \textbf{Does it run as matmul?} If the inner loop is a serial per-position recurrence with no parallel reduction, it will be 10--40$\times$ slower than the baseline and will not win. Rewrite as GEMMs or abandon.

\item \textbf{Is the behaviour designed or emergent?} If the proposal hardcodes the target phenomenon (fixed coefficients, designed operators, prescribed schedules), lower the prior. If it builds an affordance and lets the optimiser populate it, raise the prior.

\item \textbf{What is the one-epoch canary?} Specify, in advance, what can go wrong and how to detect it cheaply. If the answer is ``we'll know at the end of training,'' do not launch.

\item \textbf{Has literature solved the same problem?} Thirty minutes of reading is worth more than ten days of GPU time when the target is well-studied. The 1D-RoPE-on-patches mistake was solved publicly by Gemma, Qwen2-VL, and LLaVA before we designed V3.

\item \textbf{Is there a biological analogue?} Not as proof, but as a prior. Systems that have been under evolutionary selection for a task (memory, language, attention) sometimes encode a structural constraint that is not obvious from first principles. See Section~\ref{sec:biology}.
\end{enumerate}

Applied retrospectively, these rules would have killed most of the red rows in Section~\ref{sec:record} before they ran. Applied prospectively, they are the triage filter for future candidates.

% ======================================================================
\section{Biological Precedent}
\label{sec:biology}

Two biological frameworks recur across the productive experiments. Neither is invoked as justification ,  the architecture must stand on engineering merits ,  but both provide independent evidence that the observed principles are not accidents of our particular programme.

\subsection{Hemispheric lateralisation (McGilchrist)}

McGilchrist's claim is that intelligence requires two complementary processing modes: a \emph{literal} one (left hemisphere: categorical, sequential, detail-focused, good at prediction of known things) and a \emph{relational} one (right hemisphere: holistic, relational, context-sensitive, good at novel structure). Neither alone is sufficient; each would produce a recognisable clinical failure mode. The hemispheres communicate through a bandwidth-limited channel (corpus callosum, $\sim 200$M axons against $\sim 20$B neurons per side), and the bottleneck is argued to \emph{cause} specialisation.

The ORN programme converges on the same architecture from engineering pressure. L (the ORN proper) is the literal processor: autoregressive, token-level, trained by cross-entropy. R (the second hemisphere in \emph{Two Hemispheres}) is the relational processor: non-causal, structure-compressing, trained by coherence on L's coupling dynamics. The coupling between them is a deliberately low-rank channel ($r_{\mathrm{attn}} \ll d$); the gate settles to 0.85, not 1.0, mirroring corpus-callosum moderation.

The mapping to our rulebook:
\begin{itemize}[nosep]
\item R1 (invariance) corresponds to hemispheric specialisation: each side encodes a shared invariant of its processing mode rather than duplicating both.
\item R2 (orthogonality) corresponds to the non-overlap of representations: the right hemisphere's structural state is \emph{not} a compressed copy of left-hemisphere content; it is the orthogonal component.
\item R5 (emergence) corresponds to development: the division of labour between hemispheres matures through interaction, not through prescription.
\end{itemize}

The failure modes of our experiments also map: architectures that gave R too much bandwidth (additive injection, cross-attention) produced either L-dominance (gate closes) or content-copying (R learns L's features). The biological precedent would have predicted both outcomes.

\subsection{Complementary learning systems (McClelland, O'Reilly, Kumaran, Hassabis)}

CLS theory asserts that the brain solves the stability--plasticity dilemma by separating two learning systems: a slow, distributed neocortex that generalises across experiences, and a fast, sparse, indexed hippocampus that encodes specific episodes and replays them to the cortex for consolidation. The hippocampus compresses episodes into low-dimensional indices; during sharp-wave ripples these are replayed and absorbed by the cortex, with amplitude bounds preventing runaway interference.

The gauge-LORN v10 mechanism is structurally identical. L (frozen at inference, slow at training) is the neocortex. G (the fast learner that produces $\alpha(x)$, a low-rank gate) is the hippocampus. The unit-norm projection on $\alpha$ is the amplitude bound that prevents the new encoding from leaving the recoverable orbit ,  exactly the role of SWR amplitude regulation in the biological story.

The match is strict enough to predict the failure modes of alternative designs:
\begin{itemize}[nosep]
\item v8 (overfitting L + $\alpha$ collapse) corresponds to cortical disruption: the slow system is perturbed during fast-learning training, equivalent to the retrograde amnesia produced by consolidation failure.
\item v9 ($\|\alpha\| \to 16$, val $+1.71$ nats) corresponds to unbounded hippocampal replay overwhelming the cortex ,  a scenario CLS theory explicitly forbids.
\item v10 (unit-norm, $\|\alpha\| = 1$) corresponds to bounded replay amplitude ,  the only regime where consolidation succeeds.
\end{itemize}

The biological prediction and the empirical result align so tightly that we now treat CLS as an architectural prior: any fast-slow learning decomposition must bound the fast component's action in a way the slow component can absorb.

\subsection{Temporal binding as the brain's trick}

Perhaps the deepest biological lesson is that the brain does not use category labels. It uses \emph{temporal co-occurrence}: things that happen together in time share a hippocampal index, and categories emerge over many replay cycles. Our v11 temporal-binding experiment (Wikipedia vs OpenWebText as separate episodes, no category supervision) reproduced this empirically: it outperformed the label-supervised v10 variant, despite having strictly less information.

The lesson for R2 (orthogonality): category labels live in the NTP gradient's span almost by construction, because NTP already learns to predict category-marked tokens. Temporal co-occurrence, by contrast, is a relation between contexts, not a property of tokens, and is genuinely orthogonal to NTP. The brain appears to have arrived at this distinction by selection pressure; we arrived at it by repeatedly failing to make label supervision work.

% ======================================================================
\section{The GPU as a Selection Force}
\label{sec:gpu}

A subtle theme through the record is that \emph{the accelerator is part of the theory}. Hardware constraints are not a nuisance we route around; they actively shape which designs can be trained at all, and therefore which architectures the field converges on.

The transformer succeeded in large part because it is a stack of dense matmuls with a softmax and a GeLU. Everything about it ,  the quadratic attention that GPUs love because it is $N^2$ GEMMs; the per-layer parameter redundancy that fills tensor cores; the fused attention kernels that drop memory traffic by an order of magnitude ,  is hardware-shaped. Architectures that violate these constraints (linear attention before careful kernel work, state-space models before parallel scan, recursive or iterative operators) underperformed for years before the hardware or the algorithms caught up.

Our negative results are consistent. Every failed experiment that claimed a representational advantage but was $10\times$--$43\times$ slower to train (abelian attention, iterative field recurrence, per-token spectral composition) never got enough compute to falsify its claim. It is formally possible that one of them is correct and we simply could not afford to find out. But the correct response to that possibility is not ``train it anyway''; it is ``find the version that runs as matmul, or accept that the idea lies outside the current hardware envelope.''

The ORN is itself an example of hardware-shaped design. SharedM wins partly because it reduces parameter count but mostly because it turns $L$ separate attention projections into a single cached GEMM ,  fewer kernel launches, better tensor-core occupancy, cleaner gradient flow. The wide shared FFN wins for the same reason: one large matmul at $8\times d$ dominates the arithmetic intensity of $L$ small ones at $4\times d$. The representational argument (one wide memory vs $L$ narrow copies) and the hardware argument (one fat matmul vs $L$ thin ones) point the same way.

The rulebook implication: R4 is not a practical compromise. It is a selection principle. Architectures that run badly on GPUs are not trained enough to be tested against architectures that run well. The design space is filtered by hardware before it is filtered by loss. Any theory of what can be learned must include this filter.

% ======================================================================
\section{Putting It Together: What Is Learnable?}
\label{sec:learnable}

The six clusters, the rulebook, the biological correlates, and the GPU filter converge on a short answer to ``what is learnable in a gradient-based, GPU-parallel, next-token-predicting system?''

\begin{itemize}[nosep]
\item A quantity is learnable if it is \textbf{invariant across a dimension the model is redundant over}, and if the loss has \textbf{gradient access} to it.
\item A component is learnable if its behaviour is an \textbf{emergent fixed point} of an affordance the optimiser can populate, not a prescription of what it should do.
\item An architecture is trainable if its inner loop \textbf{decomposes into wide matmuls} the accelerator can saturate.
\item A structural axis is discoverable if the probe sits in the \textbf{orthogonal complement of the dominant training gradient}, or the NTP signal will absorb it.
\item A two-system decomposition is stable only if the fast system's action on the slow system is \textbf{amplitude-bounded} enough for the slow system to absorb by ordinary training.
\end{itemize}

These are not independent. Invariance is what allows sharing, which is what GPUs love. Orthogonality is what prevents the fast system from overwriting the slow. Emergence is how the optimiser finds the invariant without being told. The GPU filter selects for designs in which all of the above are true simultaneously. This is not a coincidence ,  it is why the transformer, the ORN, and (we predict) the two-hemisphere architecture all converge on roughly the same structural template.

% ======================================================================
\section{What We Are Now Willing to Ignore}

A theory of experiments is also a theory of \emph{what not to try}. Based on the record, we are now willing to de-prioritise an entire class of candidate research:

\begin{itemize}[nosep]
\item \textbf{Closed-form algebraic replacements for depth.} We have run five variants (MatExp, Spectral, Poly, S4-style, Deferred ORN) and all lose. The mechanism is understood (depth is re-coupling); no amount of parameter tuning will recover it.
\item \textbf{Strict-orbit gauge probes on single weight matrices.} Five variants failed with the same cosine-collapse signature. The structure does not live in individual parameters at rest; it lives in the forward pass.
\item \textbf{Designed spectral shapes.} Every attempt to prescribe $\M$'s eigenstructure has been outperformed by letting it emerge. We no longer propose experiments that fix spectrum by construction.
\item \textbf{Per-position iterative or sequential composition.} The GPU filter rules these out regardless of representational merit. Candidates must reduce to parallel reductions.
\item \textbf{Label-supervised structural probes.} Supervision from categories that NTP already discriminates is in the wrong subspace. We use temporal-binding, permutation-invariance, or magnitude-normalisation instead.
\item \textbf{Additive residual injection for structural context.} L closes the gate (0.16); the bandwidth is real but the channel is ignorable. Structural context must enter via coupling, not content.
\end{itemize}

This is, in practice, the research agenda that remains: invariance-led architectures (ORN, two-hemisphere), emergence-led diagnostics (recoverable orbits, permutation-invariant fingerprints), GPU-native operations (fused SharedM, wide shared FFN, SDPA), and two-system decompositions with bounded fast-system action (gauge-LORN v10, CLS-shaped adapters). Every other direction in the space has paid a tuition of compute and has not returned a dividend.

% ======================================================================
\section{Closing}

The goal of this paper is not a new model. It is to make explicit the priors that now shape our choice of experiments, so that future work on related problems ,  our own and others' ,  can avoid the calibration cost. The six clusters, the eight rules, the two biological frameworks, and the GPU filter are offered as a checklist. The empirical record of Section~\ref{sec:record} is offered as the dataset on which the checklist was trained.

The meta-claim is that the distinguishing property of productive neural-architecture experiments is \textbf{invariance under a dimension of redundancy, expressed orthogonally to the dominant training gradient, as an emergent fixed point of a GPU-native parameterisation}. Every cluster we have identified, including the biological ones, is a specialisation of this. Every productive experiment we have run is an instance of it. Every failed experiment missed at least one clause.

This is not a law. It is a rule of thumb with a story. But it has been strong enough, over the past year, to change what we propose before we run it ,  and that, rather than any individual result, is the most valuable thing we have taken from the programme.

\end{document}
